\subsubsection{Checkpoint Dynamics}
Figure~\ref{fig:qwen4accuracy} reveals checkpoint-dependent effects. At Step150, Block3 improves 4B Base MATH500 Avg@8 by 1.85 points (pointwise interval $[0.70,3.00]$) and AIME24 by 2.92 ($[0.42,5.83]$), but reduces AIME25 by 4.17. At Step100, the MATH500 difference is $-1.40$ ($[-2.57,-0.25]$); at Step50, only AMC23 favors Block3. The official Instruct pair also reverses: Step50 favors Token on three tasks, and Step150 favors Token on AIME24. The intervention changes the learning path, not a fixed performance offset. We retain all checkpoints rather than treating their maximum as a validated selection rule.

\subsubsection{Why Three Tokens? Block-Length Selection}
\label{sec:blockchoice}
Table~\ref{tab:blocksizes} gives the exploratory sweep that motivated retaining $k=3$. Block3 exceeds Block5 on all four Avg@8 metrics, while Block10 degenerates. Relative to Token, Block3 improves three Avg@8 metrics but decreases AIME24; Pass@8 is mixed. Thus three tokens is an empirical short-window default, not a universally optimal width.

\begin{table}[htbp]
\centering\small
\caption{Historical block-length sweep: Qwen3-1.7B-Base student, 4B-GRPO teacher, DAPO, 200 updates, seed 21. Cells show Avg@8 / Pass@8 (\%); this is not the later same-host comparison.}
\label{tab:blocksizes}
\begin{tabular}{@{}lrrrr@{}}
\toprule
Method ($k$) & MATH500 & AIME24 & AIME25 & AMC23 \\
\midrule
Token (1) & 69.12 / 88.60 & 13.33 / 26.67 & 6.67 / 16.67 & 36.75 / 62.65 \\
\textbf{Block3 (3)} & 69.27 / 87.80 & 9.58 / 30.00 & 8.75 / 26.67 & 39.91 / 61.45 \\
Block5 (5) & 67.40 / 87.60 & 9.17 / 20.00 & 7.92 / 23.33 & 34.19 / 61.45 \\
Block10 (10) & 0.00 / 0.00 & 0.00 / 0.00 & 0.00 / 0.00 & 0.00 / 0.00 \\
\bottomrule\end{tabular}

\par\smallskip\footnotesize
Token/Block3/Block5 used A800; Block10 used A100. Evaluation request seeds were not fixed explicitly. Block5 restarted with a lower rollout-memory budget; its scores survive as a curated summary. Protocol details: Appendix~\ref{history:qwen17_sweep}.
\end{table}

Larger blocks add cross-credit terms and change clipping (Equation~\ref{eq:cross}). Subsequent Block10 runs also degenerate with finite gradient maxima of 867.50 and 865.21; their built-in AMC23 grading has a caveat (Appendix~\ref{app:history}). These observations do not isolate a cause. Random3/Sliding3 test partitioning alternatives, not additional fixed widths or isolated components.

\subsubsection{Prompting and Local Feedback Compatibility}
The first 4B attempt repeats and hits the length cap before any update, excluding the block loss as the cause of those initial defects. The repair removes ChatML and separates request seeds together; it is not a chat-delimiter ablation. Appendix~\ref{app:probes} documents model-specific input probes.

Repaired 4B training-cap incidence is 263/6,400 (4.11\%) for Token versus 247/6,400 (3.86\%) for Block3. Instruct engine length stops are 94/6,400 (1.47\%) versus 77/6,400 (1.20\%), without generated thinking tags. Evaluation differences in missing boxes and length stops remain benchmark-dependent (Appendix~\ref{app:instruct}); length control is distinct from correctness.

\subsubsection{Credit Redistribution and Training Dynamics}
Across 200 updates, 4B Block3's median sign disagreement is 15.57\%, or 6.68\% with absolute-advantage weighting. Median relative redistribution is 1.148, a magnitude ratio that can exceed one; size-one Token values are zero by construction. These quantify the intervention, not feedback correctness.

Figure~\ref{fig:qwen4training} shows low late-training entropy and reduced gradients without Block10-like degeneration: redistribution alone does not diagnose collapse. Appendix~\ref{app:positionfigures} maps entropy and overlap with occupancy masks; changing on-policy prefixes preclude a fixed-input causal interpretation. Separating shared signals, clipping and loss scale requires component controls and replication (Appendix~\ref{sec:limits}).
